% 8. Conclusion (transferred from docs/paper/paper_draft.md)
\section{Conclusion}
\label{sec:08_conclusion}

Under the evaluated setting, the Drosophila visual connectome did not provide a measurable topology-specific advantage over degree-preserving rewired and other control graphs for glioma-versus-meningioma MRI classification. The central evidence is the leakage-safe five-fold cross-validation: a mean paired patient-level ROC AUC difference of −0.002 ± 0.030 against the rewired control, higher in 12 of 25 runs and lower in 13.

Graph-based models nevertheless achieved higher ROC AUC than the simple CNN baseline. This suggests that the architectural graph block may be useful independently of biological wiring.

Future work includes:

\begin{itemize}
\item the pre-declared numeric-ID subgroup analysis and the w≥5 sensitivity graph;
\item weighted and neurotransmitter-signed propagation;
\item node-type embeddings;
\item stronger encoders;
\item larger and independent MRI cohorts with external validation;
\item tasks whose structure better matches the biological circuit.
\end{itemize}
